% TOP_PROBLEM: {"id": 2302060, "title": "Research Problems in Function Theory — Problem 2.60", "category_id": 8, "category": {"id": 8, "name": "analysis", "display_name": "Analysis", "description": "Limits, continuity, calculus, and function theory.", "slug": "analysis", "order_index": 8, "created_at": "2026-07-31T15:26:25.671Z"}, "status": "open", "problem_number": "AMR-022-2060", "set_id": 13, "statusQualification": "The intended question excludes constant functions, which would otherwise give a trivial counterexample."}
% FIRST_POSTED: 2026-10-02
% ENTRY_KIND: statement_only
% UnsolvedMath ID 2302060; source catalogue code AMR-022-2060
% !TeX program = lualatex
% Definition and sources only; this file is not an LLM solution attempt.
\documentclass[11pt,a4paper]{article}
\usepackage[margin=25mm]{geometry}
\usepackage{fontspec}
\setmainfont{lmroman10-regular.otf}[BoldFont=lmroman10-bold.otf,
 ItalicFont=lmroman10-italic.otf,BoldItalicFont=lmroman10-bolditalic.otf]
\usepackage{amsmath,amssymb,mathtools}
\usepackage{unicode-math}
\setmathfont{latinmodern-math.otf}
\AtBeginDocument{\let\setminus\smallsetminus}
\usepackage{xurl,hyperref}
\hypersetup{colorlinks=true,urlcolor=blue}
\setcounter{secnumdepth}{0}
\setlength{\parindent}{0pt}
\setlength{\parskip}{5pt}
\setlength{\emergencystretch}{3em}
\begin{document}
\section*{Research Problems in Function Theory — Problem 2.60}\label{2302060}
{\small UnsolvedMath ID 2302060 · AMR-022-2060 · Analysis · AMR Open Problem Lists}
\subsection{Definitions and mathematical statement}
Let $f(z)=\sum_{k=0}^{\infty}a_kz^k$ be a nonconstant entire function with $f(z)\ne0$ for all $z\in\mathbb C$. Such a function is necessarily transcendental. For $n\ge1$ let
$S_n(z)=\sum_{k=0}^na_kz^k$ denote its Taylor section at the origin.
Must it be true that for every $\varepsilon>0$ there are an integer $n\ge1$ and a complex number $z_0$ with
\[
S_n(z_0)=0,\qquad |f(z_0)|<\varepsilon?
\]
Equivalently, the infimum of $|f(z)|$ over the union of the zero sets of all nonconstant Taylor sections should be zero. Constant sections have no zeros because $a_0=f(0)\ne0$ and can simply be ignored.
The nonconstant assumption removes the trivial constant counterexample. The function $f$ has no zeros, so the required points only approach small nonzero values of $f$. They cannot remain in a fixed compact set, since $f$ has a positive minimum modulus there. Uniform convergence of sections on compact sets therefore does not settle the problem. The source notes that if the zero-free assumption is dropped and $f$ has a zero, the corresponding assertion follows from Hurwitz's theorem.
\subsection{Short English statement}
Must some zeros of Taylor sections of a nonconstant zero-free entire function occur where the original function is arbitrarily small?
\subsection{Sources}
\begin{itemize}
\item[S1] ULAM.AI and the UnsolvedMath Contributors, supplied problems.json, record 2302060 (AMR-022-2060), AMR Open Problem Lists. Catalogue identity and source transcription; CC BY 4.0.
\url{https://huggingface.co/datasets/ulamai/UnsolvedMath}.
\item[S2] Hayman - Research Problems in Function Theory (2018), Problem 2.60. Original source indexed by AMR-022-2060.
\url{https://arxiv.org/abs/1809.07200}.
\end{itemize}
\end{document}
